\documentclass[12pt,a4paper]{article}
\usepackage[utf8]{inputenc}
\usepackage[vietnamese]{babel}
\usepackage{amsmath}
\usepackage{graphicx}
\usepackage{geometry}
\geometry{margin=2.5cm}

\title{Bai Thu Nghiem LaTeX}
\author{Hoc Vien MSA-FPT}
\date{\today}

\begin{document}

\maketitle

\section{Gioi Thieu}
Day la file LaTeX thu nghiem. Neu ban doc duoc file PDF nay, LaTeX da duoc cai dat thanh cong.

\section{Cong Thuc Toan Hoc}
Cong thuc Euler noi tieng:
\begin{equation}
    e^{i\pi} + 1 = 0
\end{equation}

Tong binh phuong:
\begin{equation}
    \sum_{i=1}^{n} x_i^2 = x_1^2 + x_2^2 + \cdots + x_n^2
\end{equation}

\section{Bang Mau}
\begin{table}[h]
\centering
\caption{Vi du bang so lieu}
\begin{tabular}{|l|c|r|}
\hline
\textbf{Phuong phap} & \textbf{Do chinh xac (\%)} & \textbf{Thoi gian (s)} \\
\hline
PP A & 78.5 & 120 \\
PP B & 85.2 & 85 \\
PP de xuat & \textbf{94.7} & \textbf{12} \\
\hline
\end{tabular}
\end{table}

\section{Ket Luan}
LaTeX da san sang de viet bao cao NCKH!

\end{document}
